\documentclass[11pt]{article}
\usepackage[margin=1in]{geometry}
\usepackage{amsmath, amssymb, amsthm}
\usepackage{algorithm}
\usepackage{algpseudocode}

\title{Problem 467: Superinteger}
\author{}
\date{}

\newtheorem{theorem}{Theorem}
\newtheorem{lemma}[theorem]{Lemma}

\begin{document}
\maketitle

\section{Problem Statement}

An integer $s$ is a \textbf{superinteger} of $n$ if the digits of $n$ form a subsequence of the digits of $s$. Let $p(i)$ be the $i$-th prime, $c(i)$ the $i$-th composite. Define the digital root $\mathrm{sd}(n) = 1 + (n-1) \bmod 9$ for $n \ge 1$, and set $\mathrm{sp}(i) = \mathrm{sd}(p(i))$, $\mathrm{sc}(i) = \mathrm{sd}(c(i))$.

Let $P_n = \mathrm{sp}(1)\cdots\mathrm{sp}(n)$ and $C_n = \mathrm{sc}(1)\cdots\mathrm{sc}(n)$. Define $f(n)$ as the numeric value of the Shortest Common Supersequence (SCS) of $P_n$ and $C_n$.

Given: $f(10) = 2357246891352679$, $f(100) \bmod (10^9+7) = 771661825$. Find $f(10000) \bmod (10^9+7)$.

\section{Mathematical Foundation}

\begin{theorem}[Digital root]
\label{thm:dr}
For $n \ge 1$, $\mathrm{sd}(n) = 1 + (n-1) \bmod 9$.
\end{theorem}

\begin{proof}
We have $\mathrm{sd}(n) \equiv n \pmod{9}$ with $\mathrm{sd}(n) \in \{1, \ldots, 9\}$. If $9 \mid n$, then $\mathrm{sd}(n) = 9$ and $(n-1) \bmod 9 = 8$, so $1 + 8 = 9$. Otherwise, $\mathrm{sd}(n) = n \bmod 9 = 1 + (n-1) \bmod 9$.
\end{proof}

\begin{theorem}[SCS--LCS duality]
\label{thm:scs}
For strings $A$ and $B$,
\[
|\mathrm{SCS}(A, B)| = |A| + |B| - |\mathrm{LCS}(A, B)|.
\]
\end{theorem}

\begin{proof}
An SCS must contain every character of $A$ and every character of $B$. Characters belonging to a common subsequence can be shared. Maximizing the shared portion (i.e., using an LCS) minimizes the total length.
\end{proof}

\begin{lemma}[LCS recurrence]
\label{lem:lcs}
For sequences $A = a_1 \cdots a_m$ and $B = b_1 \cdots b_n$:
\[
L(i,j) = \begin{cases}
0 & \text{if } i = 0 \text{ or } j = 0, \\
L(i{-}1, j{-}1) + 1 & \text{if } a_i = b_j, \\
\max(L(i{-}1,j),\, L(i,j{-}1)) & \text{otherwise}.
\end{cases}
\]
\end{lemma}

\begin{proof}
Standard DP: if $a_i = b_j$, extending the LCS of $A[1..i{-}1]$, $B[1..j{-}1]$ by the shared character is optimal. Otherwise, at least one endpoint is excluded, giving the maximum of the two subcases.
\end{proof}

\begin{theorem}[Modular evaluation]
\label{thm:mod}
The numeric value of the SCS string $d_1 d_2 \cdots d_L$ modulo $M$ is computed via Horner's method:
\[
\mathrm{val} = \bigl(\cdots((d_1 \cdot 10 + d_2) \cdot 10 + d_3) \cdots\bigr) \bmod M.
\]
\end{theorem}

\begin{proof}
This computes $\sum_{k=1}^{L} d_k \cdot 10^{L-k} \bmod M$ by the standard Horner evaluation, with all intermediate results reduced modulo~$M$.
\end{proof}

\section{Algorithm}

\begin{algorithm}[H]
\caption{Compute $f(n) \bmod M$}
\begin{algorithmic}[1]
\State Generate first $n$ primes and composites; compute $P_n$ and $C_n$ via digital roots
\State Compute LCS table $L[0..n][0..n]$ using Lemma~\ref{lem:lcs}
\State Backtrack through $L$ to construct the SCS digit string $S$
\State Compute $\mathrm{val} \gets 0$
\For{each digit $d$ in $S$}
    \State $\mathrm{val} \gets (10 \cdot \mathrm{val} + d) \bmod M$
\EndFor
\State \Return $\mathrm{val}$
\end{algorithmic}
\end{algorithm}

\section{Complexity Analysis}

\begin{itemize}
    \item \textbf{Time:} $O(n^2)$ dominated by the LCS DP table. Sieving is $O(n \log n)$; SCS construction and evaluation are $O(n)$.
    \item \textbf{Space:} $O(n^2)$ for the full DP table (needed for backtracking).
\end{itemize}

\section{Verification}

\begin{itemize}
    \item $P_{10} = 2357248152$, $C_{10} = 4689135679$
    \item $f(10) = 2357246891352679$ \checkmark
    \item $f(100) \bmod (10^9+7) = 771661825$ \checkmark
\end{itemize}

\section{Answer}

\[
\boxed{775181359}
\]

\end{document}
